\section{Artifact map and reproduction}\label{app:artifacts}
\small
The evidence snapshot is repository commit
\path{4619bb97f63d960d40ffce127154b4f54195135d}.
Published source paths below have repository prefix \texttt{crypto:}; hyperlinks resolve to
that immutable snapshot. The accompanying \texttt{data/evidence\_map.json}
records exact SHA-256 hashes and sizes, and four unmodified JSON records and one previously unpublished proof note
are copied under \texttt{data/}. No frozen input or worker was modified.

\newcommand{\artifact}[2]{\paragraph{#1}\url{https://github.com/aburan28/crypto/blob/4619bb97f63d960d40ffce127154b4f54195135d/#2}}
\paragraph{E0: actual input.}
The full chain file is
\path{crypto:research/advance/instances/n83-d18-m5-T001-offset0.json},
SHA-256
\path{5c9be468c207cc2b013123b4265ee0963e96490c16aad1b1650a3528c0ce0c25}.
The publication excluded this bulk input; a full chain replay requires
that exact file or an independently checked export. This paper includes
its parameter metadata, not its complete equations.
The norm target is the negative of the input affine point:
\[
 z=\mathtt{0x622425e96ba633969263b},\qquad
 y_T=\mathtt{0x507c02b90217be9ec635f}.
\]
The field, basis, target, and curve are recorded in
\path{data/fixture_metadata.json}; this suffices to specify the mathematical
norm predicate, without implying a byte-identical original-chain export.

\artifact{E1: binary norm chart and earlier support work}{research/full_binary_norm_circuit_20261017_v1/PROOF.md}
The adjacent \path{CONTROLS_REVIEWED.json} records the original norm oracle.
The compact support proof is copied unchanged at \path{data/support_predicate_proof.md} from the previously unpublished local \path{research/support_cost_20261016_v1/PROOF.md}; the $2^{73}$ affine-family accounting is in \path{research/actual_norm_u_rank_20261017_v1/PROOF.md}. The earlier support accounting is in
\path{research/autolab_cycle_20261017_v4/REPORT.md}; earlier divisor
and support studies are historical evidence, not newly benchmarked results.

\artifact{E2: native-coordinate profile}{lab/ledger/attempts/native-cut-aware-hybrid-norm-profile-n83-20261017-v1/0001/scientific_result.json}
Compiler proof, controls, input freezes, and source manifests are under
\path{research/hybrid_norm_xor_20261017_v1/}.

\artifact{E3: hard native search}{lab/ledger/attempts/native-hybrid-limit2-norm-search-n83-20261017-v2/0001/scientific_result.json}
The adjacent \path{process.json} gives the cold lifetime and RSS.
The driver proof and controls are under
\path{research/hybrid_norm_search_20261017_v2/}.

\artifact{E4: exact symmetry controls}{research/norm_source_symmetry_20261017_v1/PREPARATION_CONTROLS.json}
The adjacent proof and \path{SORTED_COMPARATOR_NEXT.md} separate implemented
restoration from the unrun comparator search.

\artifact{E5: wide odd-field supplied quintics}{lab/ledger/attempts/native-odd-interval-roots-NIST-panel-20261017-v2/0001/scientific_result.json}
The proof and small controls are under
\path{research/odd_interval_roots_20261017_v2/}.

\artifact{E6: preserved correction}{research/actual_norm_u_rank_20261017_v1/ERRATUM.md}
\artifact{E7: immutable release floor}{performance/GATE.md}

\paragraph{Build.}
From \path{papers/norm-rootfinding/}, run \texttt{make check}.
The build uses the shared \path{papers/paper.mk}, resolves bibliography and
cross-references, and rejects LaTeX errors and overfull boxes.
The presentation-only \path{figures/render_figures.py} reads the copied
profile JSON; it performs no field arithmetic, oracle evaluation, or timing.
Run it from the paper directory to recreate the PDF/SVG/PNG figures.
The committed visual record also appears in
\path{benchmarks/20261017-rootfinding-paper-synthesis/}.
No new solver experiment is required to reproduce these manuscript figures.

\paragraph{Native experiment replay.}
Use the source/input freezes belonging to each experiment, not a newly
compiled binary described as the old worker. The preserved scheduler entry is
\texttt{python -m lab.runner resume --ledger lab/ledger}
followed by \texttt{--manifest lab/research\_program.json}; scheduling
\texttt{run-ready} is a distinct action and can launch queued work.
The Python scheduler is legacy orchestration, not a new mathematical
implementation introduced by this manuscript. Shared-library and third-party
source availability, the original input, and the pinned native build must be
resolved before claiming a reproducible fresh hard run. Merely building this
paper does not replay those experiments.

\begin{table}[!t]
\centering\small
\begin{tabular}{p{0.2\linewidth}p{0.43\linewidth}p{0.26\linewidth}}
\toprule
Statement & Exact scope and evidence & Formal proof status\\
\midrule
Theorem~\ref{thm:support}& Finite $K,D$, monic nonconstant $H$,
$k\geq\deg H$; written multiplicity proof & No Lean declaration\\
Theorem~\ref{thm:chart}& Characteristic two, finite field, $b\ne0$,
valid target, $z\ne0$; coefficient proof and literal norm controls &
No Lean declaration\\
Theorem~\ref{thm:compiler}& Acyclic exact parity/AND definitions with
constants; complete small projections \ev{E2}& No Lean declaration\\
Proposition~\ref{prop:symmetry}& Identical domains and symmetric predicate;
full orbit restoration \ev{E4}& No Lean declaration\\
Algorithm 1 coverage & Norm chart only; solver UNSAT trusted, budgets
clear exact counts & No independent large UNSAT proof\\
\bottomrule
\end{tabular}
\caption{Proof obligations are not substituted by a weaker theorem or by
the earlier decomposition manuscript's formalization.}
\end{table}
